\section{Verificación y Validación}

\subsection{Resultados experimentales}
\label{sec:resultados}

\noindent\fbox{\parbox{\dimexpr\textwidth-2\fboxsep\relax}{%
\textbf{Hallazgo principal.} El analizador sintáctico \texttt{triton\_parser} alcanza \textbf{97.1\% PARSE\_OK} sobre el benchmark TRITONHEAL (68/70 kernels) y \textbf{100\% de éxito} en la suite unitaria (12/12). Los únicos dos rechazos corresponden a \textbf{mutaciones sintácticamente inválidas} del benchmark (\texttt{k\_0040}, \texttt{k\_0058}), confirmadas también con el analizador \texttt{ast} de Python. \textbf{No se registraron falsos rechazos} en kernels válidos ni errores léxicos en la corrida del benchmark.
}}

\vspace{0.8em}

\subsection{Resumen cuantitativo}
\begin{table}[H]
\centering
\caption{Cuadro resumen de resultados (corridas locales reproducibles con \texttt{make test} y \texttt{make benchmark}).}
\label{tab:resumen-global}
\begin{tabular}{@{}lrrp{5.2cm}@{}}
\toprule
\textbf{Experimento} & \textbf{Éxito} & \textbf{Total} & \textbf{Tasa} \\
\midrule
Suite unitaria (regresión + inválidos + léxico) & 12 & 12 & \textbf{100\%} \\
Benchmark TRITONHEAL (kernels \texttt{k\_*.tri}) & 68 & 70 & \textbf{97.1\%} \\
\quad $\hookrightarrow$ kernels válidos aceptados & 68 & 68 & \textbf{100\%} \\
\quad $\hookrightarrow$ mutaciones rechazadas (correcto) & 2 & 2 & \textbf{100\%} \\
Errores léxicos en benchmark & 0 & 70 & 0\% \\
\bottomrule
\end{tabular}
\end{table}

\subsection{Interpretación de los resultados}
Los resultados deben leerse en dos niveles:

\begin{enumerate}
  \item \textbf{Suite unitaria (12/12 PASS).} Valida que el parser cumple los entregables del curso: acepta \texttt{ejemplo\_1.tri}--\texttt{ejemplo\_3.tri} del reporte léxico, kernels Triton mínimos propios, rechaza cinco archivos con errores sintácticos deliberados, rechaza un identificador con caracteres invalidos (\texttt{BAD\_ID}, caso I6) con \texttt{LEXICAL\_ERROR}, y conserva el modo \texttt{-t} del scanner.
  \item \textbf{Benchmark de 70 kernels (68 PARSE\_OK).} Mide desempeño sobre el banco de pruebas TRITONHEAL usado en el proyecto del equipo. De 70 archivos, \textbf{68 son sintácticamente válidos} y fueron aceptados; \textbf{2 son mutaciones inválidas} sembradas en el benchmark y fueron rechazadas correctamente por el parser.
\end{enumerate}

\textbf{Conclusión:} la tasa del 97.1\% \emph{no} indica fallos en código Triton correcto; indica que el 97.1\% del \emph{total del banco} incluye 2 archivos que \emph{deben} fallar. La métrica relevante para código válido es \textbf{68/68 = 100\% de aceptación}.

\subsection{Gráfica de resultados (Figura~\ref{fig:benchmark70})}
La Figura~\ref{fig:benchmark70} presenta de forma explícita: (i) tarjetas con la tasa global y la cobertura real; (ii) desglose por categoría; (iii) diagrama de dona con el 97.1\%; (iv) mapa \emph{kernel por kernel} de los 70 archivos; (v) comparación con la suite unitaria.

\begin{figure}[H]
\centering
\includegraphics[width=\linewidth]{benchmark_parse_results.pdf}
\caption{%
\textbf{Resultados del analizador sintáctico.} Panel superior: indicadores clave. Centro izquierda: desglose (68 válidos aceptados, 2 mutaciones rechazadas, 0 falsos rechazos). Centro derecha: 97.1\% PARSE\_OK global. Inferior: mapa de los 70 kernels (\texttt{k\_0000}--\texttt{k\_0069}); verde = éxito, naranja = mutación inválida (\texttt{k\_0040}, \texttt{k\_0058}).
}
\label{fig:benchmark70}
\end{figure}

\subsection{Detalle de los 2 rechazos (comportamiento esperado)}
\begin{table}[H]
\centering
\caption{Kernels rechazados: mutaciones inválidas del benchmark.}
\label{tab:rechazos}
\begin{tabular}{@{}llp{6.5cm}@{}}
\toprule
\textbf{Kernel} & \textbf{Verificación Python} & \textbf{Motivo del rechazo} \\
\midrule
\texttt{k\_0040.tri} & \texttt{ast.parse} falla & \texttt{tl.store(out\_ptr + offs) \& ...} --- paréntesis inválido. \\
\texttt{k\_0058.tri} & \texttt{ast.parse} falla & \texttt{tl.store(...), mask=...)} --- coma mal ubicada. \\
\bottomrule
\end{tabular}
\end{table}

Estos rechazos demuestran que el parser \textbf{detecta errores sintácticos} con el mensaje \texttt{SYNTAX\_ERROR}, cumpliendo el entregable de mensajes de error del documento de evaluación.

%------------------------------------------------------------------------------
\subsection{Modelo de pruebas}
\textbf{Verificación}: ¿se construyó el producto correctamente? --- compila, ejecuta, mensajes coherentes.\\
\textbf{Validación}: ¿es el producto correcto? --- acepta Triton válido y rechaza inválido.

\subsection{Suite unitaria (\texttt{make test})}
\begin{table}[H]
\centering
\caption{Casos de prueba del equipo (11 ejecuciones).}
\begin{tabular}{@{}llc@{}}
\toprule
ID & Descripción & Resultado \\
\midrule
V1--V3 & \texttt{ejemplo\_1..3.tri} (reporte léxico) & PASS \\
V4--V5 & Kernels Triton mínimos & PASS \\
I1--I5 & Sintaxis inválida deliberada & SYNTAX\_ERROR \\
I6 & Identificador con \texttt{BAD\_ID} (\texttt{x\$}) & LEXICAL\_ERROR \\
L1 & Modo \texttt{-t} tokens & PASS \\
\midrule
\multicolumn{2}{l}{\textbf{Total PASS}} &
\IfFileExists{data/parser_pass.tex}{\input{data/parser_pass.tex}}{12} / \IfFileExists{data/parser_total.tex}{\input{data/parser_total.tex}}{12} \\
\bottomrule
\end{tabular}
\end{table}

\subsection{Evidencia: salida exitosa (\texttt{ejemplo\_1.tri})}
\IfFileExists{../results/parser/valid_ejemplo_1.tri.log}{%
\lstinputlisting[caption={Log de exito (\texttt{ejemplo\_1.tri}, extracto).}, basicstyle=\ttfamily\scriptsize, lastline=22]{../results/parser/valid_ejemplo_1.tri.log}
}{}

\subsection{Evidencia: error sintáctico (\texttt{def\_sin\_dos\_puntos.tri})}
\IfFileExists{../results/parser/invalid_def_sin_dos_puntos.tri.log}{%
\lstinputlisting[caption={Rechazo sintáctico esperado (suite unitaria).}, basicstyle=\ttfamily\scriptsize, lastline=5]{../results/parser/invalid_def_sin_dos_puntos.tri.log}
}{}

\subsection{Evidencia: error léxico \texttt{BAD\_ID} (I6)}
\IfFileExists{../results/parser/invalid_identificador_caracter_invalido.tri.log}{%
\lstinputlisting[caption={Rechazo por identificador con \texttt{\$} (regla \texttt{BAD\_ID}, alineada con Analizador-lexico).}, basicstyle=\ttfamily\scriptsize, lastline=3]{../results/parser/invalid_identificador_caracter_invalido.tri.log}
}{}

\subsection{Reproducibilidad}
\label{sub:reproducibilidad}

\noindent Todo el material para reproducir el proyecto en una \textbf{PC nueva} está en el repositorio público (sin dependencias externas de otros repos):

\begin{itemize}
  \item Repositorio: \url{https://github.com/G4ndhiV/Analizador-sint-ctico}
  \item Kernels de prueba: \path{benchmark/kernels/} (70 archivos \texttt{k\_*.tri})
  \item Fixtures léxicos: \path{tests/fixtures/lexico/}
  \item CFG documentada: \path{report/cfg/triton_cfg.bnf}
  \item Resultados de referencia: \path{results/parser/}, \path{results/benchmark/}
  \item Informe PDF: \path{deliverables/Informe_Analizador_Sintactico_Equipo11.pdf}
\end{itemize}

\subsubsection{Prerrequisitos}
\begin{table}[H]
\centering
\caption{Software necesario para reproducir compilación, pruebas, gráfica e informe.}
\label{tab:repro-deps}
\small
\begin{tabular}{@{}llp{5.8cm}@{}}
\toprule
\textbf{Herramienta} & \textbf{Obligatorio} & \textbf{Instalación típica} \\
\midrule
\texttt{git} & Sí & Clonar el repositorio \\
\texttt{cc} (GCC/Clang) & Sí & Xcode CLT (macOS) o \texttt{build-essential} (Linux) \\
Bison & Sí & \texttt{brew install bison} / \texttt{apt install bison} \\
Flex & Sí & \texttt{brew install flex} / \texttt{apt install flex} \\
Python~3 + \texttt{venv} & Sí (gráfica) & \texttt{python3 -m venv .venv}; \texttt{pip install -r requirements.txt} \\
Tectonic o \texttt{latexmk} & Solo para \texttt{make pdf} & \texttt{brew install tectonic} \\
\bottomrule
\end{tabular}
\end{table}

Verificación rápida: \texttt{make check-deps} (objetivo del \texttt{Makefile} raíz).

\subsubsection{Pasos desde cero (PC nueva)}
\begin{enumerate}
  \item \textbf{Clonar:}
\begin{verbatim}
git clone https://github.com/G4ndhiV/Analizador-sint-ctico.git
cd Analizador-sint-ctico
\end{verbatim}
  \item \textbf{Comprobar dependencias (opcional):} \texttt{make check-deps}
  \item \textbf{Compilar:} \texttt{make build} $\rightarrow$ ejecutable \texttt{build/triton\_parser}
  \item \textbf{Suite unitaria:} \texttt{make test} $\rightarrow$ \textbf{12/12 PASS}; logs en \texttt{results/parser/}
  \item \textbf{Benchmark 70 kernels:} \texttt{make benchmark} $\rightarrow$ \texttt{results/benchmark/summary.json}
  \item \textbf{Gráfica:} \texttt{make figures} $\rightarrow$ \texttt{figures/benchmark\_parse\_results.pdf}
  \item \textbf{Informe PDF:} \texttt{make pdf} $\rightarrow$ ver ruta en \path{deliverables/Informe_Analizador_Sintactico_Equipo11.pdf}
\end{enumerate}

\subsubsection{Comando único (regeneración completa)}
\begin{verbatim}
make clean build test benchmark figures pdf
\end{verbatim}

\subsubsection{Salidas verificables}
\begin{itemize}
  \item \path{results/parser/summary.json}: \texttt{"pass": 12}, \texttt{"total": 12}.
  \item \path{results/benchmark/summary.json}: \texttt{parse\_ok}: 68, \texttt{total}: 70.
  \item \path{figures/benchmark_parse_results.pdf} y \path{figures/benchmark_parse_results.png}.
  \item PDF del informe en \path{deliverables/} (Apéndice~A: CFG de 91 producciones).
\end{itemize}

\subsubsection{Consulta sin recompilar}
Basta con clonar el repositorio y abrir el PDF en \path{deliverables/} y los archivos \path{results/parser/summary.json} y \path{results/benchmark/summary.json}.

\subsubsection{Notas por sistema operativo}
\begin{itemize}
  \item \textbf{macOS:} instalar \texttt{bison} y \texttt{flex} con Homebrew; el \texttt{Makefile} enlaza la librería Flex de Homebrew.
  \item \textbf{Linux:} usar \texttt{bison} y \texttt{flex} del sistema; enlace con \texttt{-lfl}.
  \item \textbf{\texttt{BAD\_ID}:} reglas en \path{lexer/lexico_equipo_11_scanner.l} (repo léxico: \url{https://github.com/G4ndhiV/Analizador-lexico}); prueba I6 en \path{tests/invalid/identificador_caracter_invalido.tri}.
\end{itemize}

Guía extendida: \path{README.md} en la raíz del repositorio.
